%%--------------------------------------------------------------------
%1--------------------------------------------------------------------
\begin{glyph}{What (何)}{what}\label{glyph:what}
What.\\\hline
The sense of \textit{how many} can also be expressed in certain compounds.\\\hline
This 漢字 can also serve as a place-holder for the 漢字 following it.
\end{glyph}
\gindex{What}{何}\strindex{7}{何}

\begin{comps}
\comprtwocone&人亻 + 可\\\hline
\comprthreecone&Man/Person + Possible (p. \pageref{glyph:possible})\\\hline
\end{comps}

\begin{remglyph}
\hooglig{What} can that \remcom{person} \remcom{possibly} do?\\\hline
\end{remglyph}

\begin{prons}
\pronsrtwocone&---&\textbf{なん、なに (nan, nani)}\\\hline
\pronsrthreecone&カ (KA)&---\\\hline
\end{prons}

\begin{remprons}
\hooglig{What} \comkun{nanny} will use \comon{cuss} words?  \comkun{None}!\\\hline
\end{remprons}

%{\middel}
% &  \mbox{()}\\\hline
\begin{somme}
何&\textit{what?; What!}&なに \mbox{(nani)}\\\hline
何時&what + time = \textit{what time?}&なん{\middel}ジ \mbox{(nan{\middel}JI)}\\\hline
{\diff}幾何学&how many + what + study = \textit{geometry}&キ{\middel}カ{\middel}ガク \mbox{(KI{\middel}KA{\middel}GAKU)}\\\hline
何気無い&what + spirit + nil = \textit{casual, unconcerned}&なに{\middel}ゲ{\middel}ない \mbox{(nani{\middel}GE{\middel}nai)}\\\hline
何人&how many + person = \textit{how many people?}&なん{\middel}ニン \mbox{(nan{\middel}NIN)}\\\hline
何人&what + person = \textit{anyone; whoever}&なん{\middel}びと \mbox{(nan{\middel}bito)}\\\hline
何回&how many + rotate = \textit{how many times?}&なん{\middel}カイ \mbox{(nan{\middel}KAI)}\\\hline
何で&\textit{why?; what for?; how?}&なん{\middel}で \mbox{(nan{\middel}de)}\\\hline
何て&\textit{how {\ldots}!; what {\ldots}!}&なん{\middel}て \mbox{(nan{\middel}te)}\\\hline
何とか&\textit{something (or other); so-and-so}&なん{\middel}とか \mbox{(nan{\middel}toka)}\\\hline
何でも&\textit{anything; whatever (one likes); everything}&なん{\middel}でも \mbox{(nan{\middel}demo)}\\\hline
何度&what + degree = \textit{how many times?; how often?}&なん{\middel}ド \mbox{(nan{\middel}DO)}\\\hline
何故{\notcov}&what + happenstance = \textit{why?}&な{\middel}ぜ \mbox{(na{\middel}ze)}\\\hline
\end{somme}

\begin{irrr}
何時&\textit{when?; at what time?; how soon?}&い{\middel}つ \mbox{(i{\middel}tsu)}\\\hline
何時も&\textit{always; usually; every time}&い{\middel}つ{\middel}も \mbox{(i{\middel}tsu{\middel}mo)}\\\hline
何時までも&\textit{forever; for good; eternally; as long as (one likes)}&い{\middel}つ{\middel}までも \mbox{(i{\middel}tsu{\middel}mademo)}\\\hline
何時か&\textit{one day; someday; in due course; in time}&い{\middel}つ{\middel}か \mbox{(i{\middel}tsu{\middel}ka)}\\\hline
\end{irrr}
\end{multicols}

\begin{decomp}{5}
あなた&は&何者&です&か。\\\hline
あなた&は&なにもの&です&か\\\hline
you&は&what kind of person&be/is&か\\\hline
\multicolumn{5}{|r|}{\textit{What are you?}}\\\hline
\end{decomp}
\pagebreak
%%--------------------------------------------------------------------
%3--------------------------------------------------------------------
\begin{multicols}{2}
\begin{glyph}{Busy (忙)}{busy}\label{glyph:busy}
Busy.
\end{glyph}
\gindex{Busy}{忙}\strindex{6}{忙}

\begin{comps}
\comprtwocone&心忄 + 亡\\\hline
\comprthreecone&Heart/Feeling + Perish/Conceal/Dead\\\hline
\end{comps}

\begin{remglyph}
\hooglig{So busy} that my \remcom{heart} is \remcom{perishing}.\\\hline
\end{remglyph}

\begin{prons}
\pronsrtwocone&\textbf{ボウ (B\={O})}&\textbf{いそが (isoga)}\\\hline
\pronsrthreecone&---&せわ (sewa)\\\hline
\end{prons}

\begin{remprons}
\hooglig{Busy} with \comon{boring} \comkun{isogamy}.\\\hline
{\warn}\textit{Isogamy} is sexual reproduction by the fusion of similar gametes.\\\hline
\end{remprons}

%{\middel}
% &  \mbox{()}\\\hline
\begin{somme}
忙しい&\textit{busy}&いそが{\middel}しい \mbox{(isoga{\middel}shii)}\\\hline
多忙&many + busy = \textit{busy; work pressure}&タ{\middel}ボウ \mbox{(TA{\middel}B\={O})}\\\hline
お多忙中&\textit{``pardon me for interrupting while you are busy''}&お{\middel}タ{\middel}ボウ{\middel}チュウ \mbox{(o{\middel}TA{\middel}B\={O}{\middel}CH\={U})}\\\hline
大忙し&large + busy = \textit{very busy (person or thing)}&おお{\middel}いそが{\middel}し \mbox{(\={o}{\middel}isoga{\middel}shi)}\\\hline
忙しい&\textit{busy; hectic; frantic}&せわ{\middel}しい \mbox{(sewa{\middel}sh\={i})}\\\hline
忙しげ&\textit{appearing busy/restless}&せわ{\middel}しげ \mbox{(sewa{\middel}shige)}\\\hline
心忙しい&heart + busy = \textit{restless}&こころ{\middel}ぜわ{\middel}しい \mbox{(kokoro{\middel}zewa{\middel}sh\={i})}\\\hline
忙中&busy + middle = \textit{in the midst of busyness}&ボウ{\middel}チュウ \mbox{(B\={O}{\middel}CH\={U})}\\\hline
忙殺{\notcov}&busy + kill = \textit{swamped with work}&ボウ{\middel}サツ \mbox{(B\={O}{\middel}SATSU)}\\\hline
\end{somme}

\begin{decomp}{4}
多忙な&生活&を&送る。\\\hline
タボウな&セイカツ&を&おくる\\\hline
busy&living/life&を&to live one's life\\\hline
\multicolumn{4}{|r|}{\textit{I lead a busy life.}}\\\hline
\end{decomp}

\end{multicols}
%%--------------------------------------------------------------------
%4--------------------------------------------------------------------
\begin{multicols}{2}
\begin{glyph}{Slow (遅)}{slow}\label{glyph:slow}
Slow.\\\hline
Relates ideas of being \textit{late} or \textit{delayed}.
\end{glyph}
\gindex{Slow}{遅}\strindex{12}{遅}

\begin{comps}
\comprtwocone&尸 + 羊 + 辶\\\hline
\comprthreecone&Corpse/Dead Body + Sheep + Walk/Road\\\hline
\end{comps}

\begin{remglyph}
\hooglig{Slow} \remcom{corpses} of zombie \remcom{sheep} \remcom{walking on the road}.\\\hline
\end{remglyph}

\begin{prons}
\pronsrtwocone&---&\textbf{おそ、おく (oso, oku)}\\\hline
\pronsrthreecone&チ (CHI)&---\\\hline
\end{prons}

\begin{remprons}
\hooglig{Slowly} \comkun{occupying} is \comkun{awesomely} \ucomon{cheap}.\\\hline
\end{remprons}

%{\middel}
% &  \mbox{()}\\\hline
\begin{somme}
遅い&\textit{slow; delayed}&おそ{\middel}い \mbox{(oso{\middel}i)}\\\hline
遅れる&\textit{be late/delayed}&おく{\middel}れる \mbox{(oku{\middel}reru)}\\\hline
出遅れる&exit + slow = \textit{to get a late start}&で{\middel}おく{\middel}れる \mbox{de{\middel}oku{\middel}reru}\\\hline
遅出&slow + exit = \textit{late attendance}&おそ{\middel}で \mbox{(oso{\middel}de)}\\\hline
手遅れ&hand + late = \textit{being too late}&て{\middel}おく{\middel}れ \mbox{(te{\middel}oku{\middel}re)}\\\hline
遅配&slow + distribute = \textit{delay in rationing}&チ{\middel}ハイ \mbox{(CHI{\middel}HAI)}\\\hline
遅かれ早かれ&slow + early = \textit{sooner or later}&おそ{\middel}かれ{\middel}はや{\middel}かれ \mbox{(oso{\middel}kare{\middel}haya{\middel}kare)}\\\hline
遅速&slow + fast = \textit{speed; progress}&チ{\middel}ソク \mbox{(CHI{\middel}SOKU)}\\\hline
遅達&slow + fee = \textit{late delivery / notification}&チ{\middel}タツ \mbox{(CHI{\middel}TATSU)}\\\hline
\end{somme}

\begin{decomp}{3}
何故&遅れた&の？\\\hline
なぜ&おくれた&の\\\hline
why&to be late&の\\\hline
\multicolumn{3}{|r|}{\textit{Why were you late?}}\\\hline
\end{decomp}

\end{multicols}

\pagebreak
%%--------------------------------------------------------------------
%5--------------------------------------------------------------------

\begin{multicols}{2}
\begin{glyph}{Fee (料)}{fee}\label{glyph:fee}
Fee.  Charge.\\\hline
A secondary term relates to \remcom{materials}.\\\hline
Ensure that you can clearly distinguish between 料 and 科 (\textit{Course}, Entry \ref{glyph:course}, p. \pageref{glyph:course}).
\end{glyph}
\gindex{Fee}{料}\strindex{10}{料}

\begin{comps}
\comprtwocone&米 + 斗\\\hline
\comprthreecone&Rice + Dipper\\\hline
\end{comps}

\begin{remglyph}
\hooglig{The fee} for \remcom{rice} depends on its \remcom{volume}.\\\hline
\end{remglyph}

\begin{prons}
\pronsrtwocone&\textbf{リョウ (RY\={O})}&---\\\hline
\pronsrthreecone&---&---\\\hline
\end{prons}

\begin{remprons}
\hooglig{The fee} of a \comon{Ry\={o}bi} watergun.\\\hline
\end{remprons}

%{\middel}
% &  \mbox{()}\\\hline
\begin{somme}
料金&fee + gold (money) = \textit{fee; charge}&リョウ{\middel}キン \mbox{(RY\={O}{\middel}KIN)}\\\hline
料理&fee + reason = \textit{cooking; cuisine}&リョウ{\middel}リ \mbox{(RY\={O}{\middel}RI)}\\\hline
入場料&enter + place + fee = \textit{admission fee}&ニュウ{\middel}ジョウ{\middel}リョウ \mbox{(NY\={U}{\middel}J\={O}{\middel}RY\={O})}\\\hline
食料品&food + goods = \textit{groceries; foodstuff}&ショクリョウヒン \mbox{(SHOKU{\middel}RY\={O}{\middel}HIN)}\\\hline
無料&nil + fee = \textit{free (of charge); gratuitous}&ム{\middel}リョウ \mbox{(MU{\middel}RY\={O})}\\\hline
有料&have + fee = \textit{not free}&ユウ{\middel}リョウ \mbox{(Y\={U}{\middel}RY\={O})}\\\hline
送料&send + fee = \textit{postage; carriage}&ソウ{\middel}リョウ \mbox{(S\={O}{\middel}RY\={O})}\\\hline
衣料&clothes + fee = \textit{clothing}&イ{\middel}リョウ \mbox{(I{\middel}RY\={O})}\\\hline
\end{somme}
\end{multicols}

\begin{decomp}{7}
急行&の&料金&は&幾ら&です&か。\\\hline
キュウコウ&の&リョウキン&は&いくら&です&か\\\hline
express&の&fee&は&how much&be/is&か\\\hline
\multicolumn{7}{|r|}{\textit{How much is the express?}}\\\hline
\end{decomp}

%%--------------------------------------------------------------------
%6--------------------------------------------------------------------
\begin{multicols}{2}
\begin{glyph}{Business (業)}{business}\label{glyph:business}
Business.  Industry.
\end{glyph}
\gindex{Business}{業}\strindex{13}{業}

\begin{comps}
\comprtwocone&並 + 䒑 + 未\\\hline
\comprthreecone&(Almost) Lined Up + Grass + Not Yet\\\hline
\end{comps}

\begin{remglyph}
\hooglig{Businesses} for \remcom{lined up} \remcom{grass} has \remcom{not yet} launched.\\\hline
\end{remglyph}

\begin{prons}
\pronsrtwocone&\textbf{ギョウ (GY\={O})}&---\\\hline
\pronsrthreecone&ゴウ (G\={O})&わざ (waza)\\\hline
\end{prons}

\begin{remprons}
\hooglig{Business} can be \ucomon{gory} when lots of people in the streets yell, ``\ucomkun{Wazzup}, \comon{get your} product today.''\\\hline
\end{remprons}

%{\middel}
% &  \mbox{()}\\\hline
\begin{somme}
工業&craft + business = \textit{industry}&コウ{\middel}ギョウ \mbox{(K\={O}{\middel}GY\={O})}\\\hline
事業&matter + business = \textit{project; venture}&ジ{\middel}ギョウ \mbox{(JI{\middel}GY\={O})}\\\hline
商業&commerce + business = \textit{commerce; business}&ショウ{\middel}ギョウ \mbox{(SH\={O}{\middel}GY\={O})}\\\hline
神業&god + business = \textit{divine work; miracle}&かみ{\middel}わざ \mbox{(kami{\middel}waza)}\\\hline
失業&lose + business = \textit{unemployment}&シツ{\middel}ギョウ \mbox{(SHITSU{\middel}GY\={O})}\\\hline
実業家&industry/business + house = \textit{businessman; industrialist}&ジツギョウカ \mbox{(JITSU{\middel}GY\={O}{\middel}KA)}\\\hline
開業&open + business = \textit{opening a business/practice}&カイ{\middel}ギョウ \mbox{(KAI{\middel}GY\={O})}\\\hline
学業&study + business = \textit{studies; schoolwork}&ガク{\middel}ギョウ \mbox{(GAKU{\middel}GY\={O})}\\\hline
\end{somme}

\begin{decomp}{4}
日本&は&工業国&だ。\\\hline
ニッポン&は&コウギョウコク&だ\\\hline
Japan&は&industrialised nation&be/is\\\hline
\multicolumn{4}{|r|}{\textit{Japan is an industrial nation.}}\\\hline
\end{decomp}

\end{multicols}
\pagebreak
%%--------------------------------------------------------------------
%7--------------------------------------------------------------------

\begin{multicols}{2}
\begin{glyph}{Swim (泳)}{swim}\label{glyph:swim}
Swim.
\end{glyph}
\gindex{Swim}{泳}\strindex{8}{泳}

\begin{comps}
\comprtwocone&水氵氺 + 永\\\hline
\comprthreecone&Water + Eternal (p. \pageref{glyph:eternal})\\\hline
\end{comps}

\begin{remglyph}
\hooglig{Swimming} in the \remcom{water} for all \remcom{eternity}.\\\hline
\end{remglyph}

\begin{prons}
\pronsrtwocone&\textbf{エイ (EI)}&\textbf{およ (oyo)}\\\hline
\pronsrthreecone&---&---\\\hline
\end{prons}

\begin{remprons}
\hooglig{Swimming} \comon{aims} in \comkun{Oyo}.\\\hline
{\warn}\textit{Oyo} is an ancient city in Nigeria.\\\hline
\end{remprons}

%{\middel}
% &  \mbox{()}\\\hline
\begin{somme}
泳ぐ&\textit{to swim}&およ{\middel}ぐ \mbox{(oyo{\middel}gu)}\\\hline
平泳ぎ&flat + swim = \textit{breaststroke}&ひら{\middel}およ{\middel}ぎ \mbox{(hira{\middel}oyo{\middel}gi)}\\\hline
背泳ぎ&back + swim = \textit{backstroke}&せ{\middel}およ{\middel}ぎ \mbox{(se{\middel}oyo{\middel}gi)}\\\hline
泳法&swim + law = \textit{swimming style}&エイ{\middel}ホウ \mbox{(EI{\middel}H\={O})}\\\hline
水泳&water + swim = \textit{swimming}&スイ{\middel}エイ \mbox{(SUI{\middel}EI)}\\\hline
遠泳&far + swim = \textit{long-distance swimming}&エン{\middel}エイ \mbox{(EN{\middel}EI)}\\\hline
遊泳&play + swim = \textit{swimming; bathing}&ユウ{\middel}エイ \mbox{(Y\={U}{\middel}EI)}\\\hline
泳動&swim + move = \textit{migration; movement}&エイ{\middel}ドウ \mbox{(EI{\middel}D\={O})}\\\hline
水泳着&swimming + wear = \textit{swimming/bathing suit}&スイ{\middel}エイ{\middel}ぎ　\mbox{(SUI{\middel}EI{\middel}gi)}\\\hline
\end{somme}

\begin{decomp}{5}
水泳&は&良い&運動&だ。\\\hline
スイエイ&は&よい&ウンドウ&だ\\\hline
swimming&は&good&exercise&be/is\\\hline
\multicolumn{5}{|r|}{\textit{Swimming is good exercise.}}\\\hline
\end{decomp}

\end{multicols}
%%--------------------------------------------------------------------
%8--------------------------------------------------------------------
\begin{multicols}{2}
\begin{glyph}{Connection (係)}{connection}\label{glyph:connection}
Connection, often with the sense of being concerned with or involved in something.\\\hline
When used on its own or as a suffix, however, this 漢字 usually indicates a person working at a certain type of job---ticket collectors, tellers, receptionists and sales clerks are a few examples of compounds making use of this character.
\end{glyph}
\gindex{Connection}{係}\strindex{9}{係}

\begin{comps}
\comprtwocone & 人亻 + 系 \\\hline
\comprthreecone & Man/Person + System (\ref{glyph:system}) \\\hline
\end{comps}

\begin{remglyph}
\hooglig{Connecting} to a \remcom{system} of \remcom{people}.\\\hline
\end{remglyph}

\begin{prons}
\pronsrtwocone & \textbf{ケイ (KEI)} & \textbf{かかり (kakari)}\\\hline
\rye{3}{かかり is the reading used as a suffix to denote the types of jobs mentioned above.  It is always voiced when in the final position (as in the second example below).}
\pronsrthreecone & --- & --- \\\hline
\end{prons}

\begin{remprons}
\hooglig{Connection} that \comon{caters} for the supply of \comkun{kakari} rugby boots.\\\hline
\end{remprons}

%{\middel}
% &  \mbox{()}\\\hline
\begin{somme}
係 & \textit{clerk} & かかり \mbox{(kakari)}\\\hline
受付係 & receive + attach + connection = \textit{receptionist; information clerk} & うけ{\middel}つけ{\middel}がかり \mbox{(uke{\middel}tsuke{\middel}gakari)}\\\hline
\notyet{関係} & checkpoint + connection = \textit{connection; relationship} & カン{\middel}ケイ \mbox{(KAN{\middel}KEI)}\\\hline
\end{somme}

\begin{decomp}{6}
無関係&の&人&を&殺る&な！\\\hline
ムカンケイ&の&ひと&を&やる&な\\\hline
unrelated&の&person&を&to do someone in&な\\\hline
\multicolumn{6}{|r|}{\textit{Don't kill off bystanders!}}\\\hline
\end{decomp}

\end{multicols}

\pagebreak
%%--------------------------------------------------------------------
%9--------------------------------------------------------------------

\begin{multicols}{2}
\begin{glyph}{Drink (飲)}{drink}\label{glyph:drink}
Drink.
\end{glyph}
\gindex{Drink}{飲}\strindex{12}{飲}

\begin{comps}
\comprtwocone & 食 + 欠 \\\hline
\comprthreecone & Eat/Food + Lack \\\hline
\end{comps}

\begin{remglyph}
\hooglig{Drinking} involves a \remcom{lack} of \remcom{sustenance}.\\\hline
\end{remglyph}

\begin{prons}
\pronsrtwocone & \textbf{イン (IN)} & \textbf{の (no)}\\\hline
\pronsrthreecone & --- & --- \\\hline
\end{prons}

\begin{remprons}
\hooglig{Drinking} \comon{ink} is for the \comkun{nobs}.\\\hline
\end{remprons}

%{\middel}
% &  \mbox{()}\\\hline
\begin{somme}
飲む & \textit{drink} & の{\middel}む \mbox{(no{\middel}mu)}\\\hline
飲み物 & drink + thing = \textit{a drink} & の{\middel}み　もの \mbox{(no{\middel}mi mono)}\\\hline
飲食 & drink + eat = \textit{drinking and eating; food and drink} & イン{\middel}ショク \mbox{(IN{\middel}SHOKU)}\\\hline
\end{somme}

\begin{decomp}{3}
何か&飲み物&でも？\\\hline
なにか&のみもの&でも\\\hline
something&drink&でも\\\hline
\multicolumn{3}{|r|}{\textit{Would you care for something to drink?}}\\\hline
\end{decomp}

\end{multicols}

%%--------------------------------------------------------------------
%10-------------------------------------------------------------------
\begin{multicols}{2}
\begin{glyph}{Yellow (黄)}{yellow}\label{glyph:yellow}
Yellow.
\end{glyph}
\gindex{Yellow}{黄}\strindex{11}{黄}

\begin{comps}
\comprtwocone & 黃 \\\hline
\comprthreecone & Yellow \\\hline
\end{comps}

\begin{remglyph}
Part of the 漢字 radicals (\#{201}).\\\hline
\end{remglyph}

\begin{prons}
\pronsrtwocone & \textbf{オウ、コウ (\={O}, K\={O})} & \textbf{き (ki)}\\\hline
\rye{3}{This is an odd character when it comes to pronunciations, as it has three common readings despite there being very few compounds containing the 漢字.  In addition, several of these compounds can be read with either of the onyomi; in such cases, though, one of them will invariable be more widely used.}
\pronsrthreecone & --- & こ (ko) \\\hline
\rye{3}{You will only encounter this reading as an alternative to the second compound below: 黄金 (こ{\middel}がね/ko{\middel}gane) \textit{gold} or \textit{money}.  Nevertheless, it is far more common for 黄金色 (こ{\middel}がね{\middel}いろ/ko{\middel}gane{\middel}iro) \textit{gold color}, to be read with the 訓読み of these 漢字.}
\end{prons}

\begin{remprons}
\hooglig{The yellow} \ucomkun{cobble-stone} road \comkun{keepers} \comon{cause} \comon{awesome} events.\\\hline
\end{remprons}

%{\middel}
% &  \mbox{()}\\\hline
\begin{somme}
黄色 & \textit{yellow} & き{\middel}いろ \mbox{(ki{\middel}iro)}\\\hline
黄金 & yellow + gold = \textit{gold; money} & オウ{\middel}ゴン \mbox{(\={O}{\middel}GON)}\\\hline
黄海 & yellow + sea = \textit{the Yellow Sea} & コウ{\middel}カイ \mbox{(K\={O}{\middel}KAI)}\\\hline
\end{somme}

\begin{decomp}{5}
その&木の葉&は&黄色く&なった。\\\hline
その&このは&は&きいろく&なった\\\hline
that / the&leaves of trees&は&yellow&to become\\\hline
\multicolumn{5}{|r|}{\textit{The leaves of the tree turned yellow.}}\\\hline
\end{decomp}

\end{multicols}

\pagebreak
%%--------------------------------------------------------------------
%11-------------------------------------------------------------------

\begin{multicols}{2}
\begin{glyph}{Answer (答)}{answer}\label{glyph:answer}
Answer.
\end{glyph}
\gindex{Answer}{答}\strindex{12}{答}

\begin{comps}
\comprtwocone & 竹ケ + 合 \\\hline
\comprthreecone & Bamboo + Match (\ref{glyph:match}) \\\hline
\end{comps}

\begin{remglyph}
\hooglig{Answers} for the \remcom{bamboo} \remcom{match}.\\\hline
\end{remglyph}

\begin{prons}
\pronsrtwocone & \textbf{トウ (T\={O})} & \textbf{こた (kota)}\\\hline
\pronsrthreecone & --- & --- \\\hline
\end{prons}

\begin{remprons}
\hooglig{Answers} are in the \comon{torture} \comkun{cottage}.\\\hline
\end{remprons}

%{\middel}
% &  \mbox{()}\\\hline
\begin{somme}
\rye{3}{The final compound below is the only instance in which the 音読み becomes voiced.}
答え & \textit{answer; response} & こた{\middel}え \mbox{(kota{\middel}e)}\\\hline
答える & \textit{to answer; respond} & こた{\middel}える \mbox{(kota{\middel}eru)}\\\hline
問答 & question + answer = \textit{questions and answers} & モン{\middel}ドウ \mbox{(MON{\middel}D\={O})}\\\hline
\end{somme}

\begin{decomp}{4}
質問&に&答え&なさい。\\\hline
シツモン&に&こたえ&なさい\\\hline
question&に&answer&to do\\\hline
\multicolumn{4}{|r|}{\textit{Answer the question.}}\\\hline
\end{decomp}

\end{multicols}

%%--------------------------------------------------------------------
%12-------------------------------------------------------------------
\begin{multicols}{2}
\begin{glyph}{Travel (旅)}{travel}\label{glyph:travel}
Travel.  Journey.
\end{glyph}
\gindex{Travel}{旅}\strindex{10}{旅}

\begin{comps}
\comprtwocone & 方 + 𠂉 + 衣\\\hline
\comprthreecone & Way + Person + Clothes\\\hline
\end{comps}

\begin{remglyph}
\hooglig{Travel} that \remcom{way} with the \remcom{person's} \remcom{clothes}.\\\hline
\end{remglyph}

\begin{prons}
\pronsrtwocone & \textbf{リョ (RYO)} & \textbf{たび (tabi)}\\\hline
\pronsrthreecone & --- & --- \\\hline
\end{prons}

\begin{remprons}
\hooglig{Traveling} to \comon{Lyon} has made me a \comkun{tubby} person.\\\hline
\end{remprons}

%{\middel}
% &  \mbox{()}\\\hline
\begin{somme}
旅 & \textit{travel; trip} & たび \mbox{(tabi)}\\\hline
旅人 & travel + person = \textit{traveler} & たび{\middel}びと \mbox{(tabi{\middel}bito)}\\\hline
旅行 & travel + go = \textit{travel; trip} & リョ{\middel}コウ \mbox{(RYO{\middel}K\={O})}\\\hline
\end{somme}

\begin{decomp}{3}
旅行&は&楽しい。\\\hline
リョコウ&は&たのしい\\\hline
travel&は&enjoyable / fun\\\hline
\multicolumn{3}{|r|}{\textit{It's fun to travel.}}\\\hline
\end{decomp}

\end{multicols}

\pagebreak
%%--------------------------------------------------------------------
%13-------------------------------------------------------------------

\begin{multicols}{2}
\begin{glyph}{Dish (皿)}{dish}\label{glyph:dish}
Dish.  Plate.
\end{glyph}
\gindex{Dish}{皿}\strindex{5}{皿}

\begin{comps}
\comprtwocone&皿 \\\hline
\comprthreecone&Dish/Plate \\\hline
\end{comps}

\begin{remglyph}
Part of the 漢字 radicals (\#{108}).\\\hline
\end{remglyph}

\begin{prons}
\pronsrtwocone & --- & \textbf{さら (sara)}\\\hline
\pronsrthreecone & --- & --- \\\hline
\end{prons}

\begin{remprons}
\hooglig{Dish} up a \comkun{sarabande}.\\\hline
\end{remprons}

%{\middel}
% &  \mbox{()}\\\hline
\begin{somme}
皿&\textit{dish}&さら \mbox{(sara)}\\\hline
皿洗い&dish + wash = \textit{dishwashing}&さら{\middel}あら{\middel}い \mbox{(sara{\middel}ara{\middel}i)}\\\hline
小皿&small + dish = small dish&こ{\middel}ざら \mbox{(ko{\middel}zara)}\\\hline
木皿&wood + dish = \textit{wooden plate}&き{\middel}ざら \mbox{(ki{\middel}zara)}\\\hline
\end{somme}

\begin{decomp}{3}
皿洗い&を&遣ります。\\\hline
さらあらい&を&やります。\\\hline
dishwashing&を&to do\\\hline
\multicolumn{3}{|r|}{\textit{I'll wash the dishes.}}\\\hline
\end{decomp}
\end{multicols}
%%--------------------------------------------------------------------
%14-------------------------------------------------------------------
\begin{multicols}{2}
\begin{glyph}{Govern (治)}{govern}\label{glyph:govern}
This interesting 漢字 conveys a sense of governing or having things under control.\\\hline
It is a meaning that can extend to managing one's health, or having an illness cured.\\\hline
Make sure to clearly distinguish this character from 法 (``law'', Entry \ref{glyph:law}), as the two are visually similar.
\end{glyph}
\gindex{Govern}{治}\strindex{8}{治}

\begin{comps}
\comprtwocone & 水氵氺 + 台 \\\hline
\comprthreecone & Water + Platform (\ref{glyph:platform}) \\\hline
\end{comps}

\begin{remglyph}
\hooglig{Govern} from the sacred \remcom{water} \remcom{platform}.\\\hline
\end{remglyph}

\begin{prons}
\pronsrtwocone & \textbf{チ、ジ (CHI, JI)} & \textbf{おさ、なお (osa, nao)}\\\hline
\rye{3}{Though we are dealing with four common readings here, things are much simpler than they appear at first glance.  The 訓読み are only found in the two intransitive/transitive verb pairs below, and チ is the more common of the 音読み.}
\pronsrthreecone & --- & --- \\\hline
\end{prons}

\begin{remprons}
\hooglig{Govern} a \comon{cheap} \comon{jig} for \comkun{Osama} and \comkun{Naomi}.\\\hline
\end{remprons}

%{\middel}
% &  \mbox{()}\\\hline
\begin{somme}
治まる (intr.) & \textit{to be tranquil; calm down} & おさ{\middel}まる \mbox{(osa{\middel}maru)}\\\hline
治める (tr.) & \textit{to govern / administer} & おさ{\middel}める \mbox{(osa{\middel}meru)}\\\hline
治る (intr.) & \textit{to get well / recover} & なお{\middel}る \mbox{(nao{\middel}ru)}\\\hline
治す (tr.) & \textit{to heal} & なお{\middel}す \mbox{(nao{\middel}su)}\\\hline
自治 & self + govern = \textit{self-government; autonomy} & ジ{\middel}チ \mbox{(JI{\middel}CHI)}\\\hline
明治 & bright + govern = \textit{Meiji (period / emperor)} & メイ{\middel}ジ \mbox{(MEI{\middel}JI)}\\\hline
\end{somme}

\begin{decomp}{3}
風&が&治まった。\\\hline
かぜ&が&おさまった\\\hline
wind&が&to lessen\\\hline
\multicolumn{3}{|r|}{\textit{The wind has calmed down.}}\\\hline
\end{decomp}

\end{multicols}

\pagebreak
%%--------------------------------------------------------------------
%15-------------------------------------------------------------------

\begin{multicols}{2}
\begin{glyph}{Confer (与)}{confer}\label{glyph:confer}
Confer.  Give.
\end{glyph}
\gindex{Confer}{与}\strindex{3}{与}

\begin{comps}
\comprtwocone & 丂 + 一 \\\hline
\comprthreecone & Reach + One\\\hline
\end{comps}

\begin{remglyph}
\hooglig{Conferred} to you because you \remcom{reached} for \remcom{one}.\\\hline
\end{remglyph}

\begin{prons}
\pronsrtwocone & \textbf{ヨ (YO)} & \textbf{あた (ata)}\\\hline
\pronsrthreecone & --- & --- \\\hline
\end{prons}

\begin{remprons}
\hooglig{Confer} \comkun{ataxia} to the \comon{yob}.\\\hline
{\warn}\textit{Ataxia} is a degenerative disease of the nervous system.\\\hline
\end{remprons}

%{\middel}
% &  \mbox{()}\\\hline
\begin{somme}
与える & \textit{confer; give} & あた{\middel}える \mbox{(ata{\middel}eru)}\\\hline
\notyet{授与} & grant + confer = \textit{conferment; awarding} & ジュ{\middel}ヨ \mbox{(JU{\middel}YO)}\\\hline
\end{somme}

\begin{decomp}{7}
私&は&犬&に&肉&を&与えた。\\\hline
わたし&は&いぬ&に&にく&を&あたえた\\\hline
I&は&dog&に&meat&を&to give\\\hline
\multicolumn{7}{|r|}{\textit{I fed some meat to my dog.}}\\\hline
\end{decomp}

\end{multicols}

%%--------------------------------------------------------------------
%16-------------------------------------------------------------------
\begin{multicols}{2}
\begin{glyph}{Daytime (昼)}{daytime}\label{glyph:daytime}
Daytime.  Midday.
\end{glyph}
\gindex{Daytime}{昼}\strindex{9}{昼}

\begin{comps}
\comprtwocone & 尺　＋　旦 \\\hline
\comprthreecone & Measurement (\ref{glyph:measurement}) + Dawn (\ref{glyph:dawn}) \\\hline
\end{comps}

\begin{remglyph}
\hooglig{Daytime} \remcom{measurements} come after \remcom{dawn}.\\\hline
\end{remglyph}

\begin{prons}
\pronsrtwocone & \textbf{チュウ (CH\={U})} & \textbf{ひる (hiru)}\\\hline
\pronsrthreecone & --- & --- \\\hline
\end{prons}

\begin{remprons}
\hooglig{Middays}, \comkun{he rules} who gets \comon{chewed} out.\\\hline
\end{remprons}

%{\middel}
% &  \mbox{()}\\\hline
\begin{somme}
昼 & \textit{daytime; midday} & ひる \mbox{(hiru)}\\\hline
昼間 & daytime + interval = \textit{daytime; during the day} & ひる{\middel}ま \mbox{(hiru{\middel}ma)}\\\hline
昼食 & daytime + eat = \textit{lunch} & チュウ{\middel}ショク \mbox{(CH\={U}{\middel}SHOKU)}\\\hline
\end{somme}

\begin{decomp}{3}
昼食&を&食べよう。\\\hline
チュウショク&を&たべよう\\\hline
lunch&を&to eat\\\hline
\multicolumn{3}{|r|}{\textit{Let's eat lunch.}}\\\hline
\end{decomp}

\end{multicols}

\pagebreak
%%--------------------------------------------------------------------
%17-------------------------------------------------------------------

\begin{multicols}{2}
\begin{glyph}{Produce (産)}{produce}\label{glyph:produce}
Produce.  Give birth to.
\end{glyph}
\gindex{Produce}{産}\strindex{11}{産}

\begin{comps}
\comprtwocone & 立 + 厂 + 生 \\\hline
\comprthreecone & Stand/Erect + Cliff + Life \\\hline
\end{comps}

\begin{remglyph}
\hooglig{Produce} new \remcom{life} while \remcom{standing} on a \remcom{cliff}.\\\hline
\end{remglyph}

\begin{prons}
\pronsrtwocone & \textbf{サン (SAN)} & \textbf{う (u)}\\\hline
\rye{3}{The 訓読み forms the intransitive/transitive verb pair 産まれる/産む (うまれる/うむ) \textit{to be born/to give birth to}; these verbs have the same meanings as the identically pronounced 生まれる/生む, which we learned back in Entry \ref{glyph:life}.\\\hline
The 音読み can become voiced in the second position, but only in a few compounds related to childbirth.}
\pronsrthreecone & --- & うぶ (ubu) \\\hline
\end{prons}

\begin{remprons}
\hooglig{Producing} \comon{sons} is the \comkun{oomph} in \ucomkun{ubuntu}.\\\hline
\end{remprons}

%{\middel}
% &  \mbox{()}\\\hline
\begin{somme}
産院 & produce + institution = \textit{maternity hospital} & サン{\middel}イン \mbox{(SAN{\middel}IN)}\\\hline
生産 & life + produce = \textit{production} & セイ{\middel}サン \mbox{(SEI{\middel}SAN)}\\\hline
水産 & water + produce = \textit{marine products} & スイ{\middel}サン \mbox{(SUI{\middel}SAN)}\\\hline
\end{somme}

\begin{irrr}
土産 & earth + produce = \textit{souvenir} & みやげ \mbox{(miyage)}\\\hline
\rye{3}{The honorary prefix お is often attached to this word:  \mbox{お土産}.}
\end{irrr}

\begin{decomp}{5}
此れ&は&お&土産&です。\\\hline
これ&は&お&みやげ&です\\\hline
this&は&honorific prefix&souvenir&be / is\\\hline
\multicolumn{5}{|r|}{\textit{This is a little gift for you.}}\\\hline
\end{decomp}

\end{multicols}

%%--------------------------------------------------------------------
%18-------------------------------------------------------------------
\begin{multicols}{2}
\begin{glyph}{Construct (建)}{construct}\label{glyph:construct}
Construct.  Build.
\end{glyph}
\gindex{Construct}{建}\strindex{9}{建}

\begin{comps}
\comprtwocone & 聿 + 廴 \\\hline
\comprthreecone & Brush + Long Stride \\\hline
\end{comps}

\begin{remglyph}
\hooglig{In constructing} something, your \remcom{inkbrush} has to walk \remcom{long strides}.\\\hline
\end{remglyph}

\begin{prons}
\pronsrtwocone & \textbf{ケン (KEN)} & \textbf{た (ta)}\\\hline
\pronsrthreecone & コン (KON) & --- \\\hline
\end{prons}

\begin{remprons}
\hooglig{Construction} of \comkun{tunnel} \comon{kennels} by a \comon{con} artist.\\\hline
\end{remprons}

%{\middel}
% &  \mbox{()}\\\hline
\begin{somme}
\rye{3}{As the second example shows, this is another 漢字 whose verb form can incorporate its ひらがな ending in certain compounds.}
建てる & \textit{construct; build} & た{\middel}てる \mbox{(ta{\middel}teru)}\\\hline
建物 & contruct + thing = \textit{building} & たて{\middel}もの \mbox{(tate{\middel}mono)}\\\hline
建国 & construct + country = \textit{founding of a country} & ケン{\middel}コク \mbox{(KEN{\middel}KOKU)}\\\hline
\end{somme}

\begin{decomp}{6}
あの&建物&は&何&です&か。\\\hline
あの&たてもの&は&なに&です&か\\\hline
that&building&は&what&be / is&か\\\hline
\multicolumn{6}{|r|}{\textit{What's that building?}}\\\hline
\end{decomp}

\end{multicols}

\pagebreak
%%--------------------------------------------------------------------
%19-------------------------------------------------------------------

\begin{multicols}{2}
\begin{glyph}{Younger Sister (妹)}{younger_sister}\label{glyph:younger_sister}
Younger sister.
\end{glyph}
\gindex{Younger Sister}{妹}\strindex{8}{妹}

\begin{comps}
\comprtwocone & 女 + 未\\\hline
\comprthreecone & Woman + Not Yet (\ref{glyph:not_yet})\\\hline
\end{comps}

\begin{remglyph}
\hooglig{My younger sister} is \remcom{not yet} a \remcom{woman}.\\\hline
\end{remglyph}

\begin{prons}
\pronsrtwocone & \textbf{マイ (MAI)} & \textbf{いもうと (im\={o}to)}\\\hline
\pronsrthreecone & --- & --- \\\hline
\end{prons}

\begin{remprons}
\hooglig{Younger sister} \comon{migrated} in her quest for \comkun{immortality}.\\\hline
\end{remprons}

%{\middel}
% &  \mbox{()}\\\hline
\begin{somme}
妹 & \textit{younger sister} & いもうと \mbox{(im\={o}to)}\\\hline
姉妹 & elder sister + younger sister = \textit{sisters} & シ{\middel}マイ \mbox{(SHI{\middel}MAI)}\\\hline
\end{somme}

\begin{decomp}{4}
妹&さん&は&元気？\\\hline
いもうと&さん&は&ゲンキ\\\hline
younger sister&\textit{polite}&は&well / fit\\\hline
\multicolumn{4}{|r|}{\textit{How's your sister?}}\\\hline
\end{decomp}

\end{multicols}

%%--------------------------------------------------------------------
%20-------------------------------------------------------------------
\begin{multicols}{2}
\begin{glyph}{Pond (池)}{pond}\label{glyph:pond}
Pond.  Reservoir.
\end{glyph}
\gindex{Pond}{池}\strindex{6}{池}

\begin{comps}
\comprtwocone & 水氵氺 + 也 \\\hline
\comprthreecone & Water + To be classical/Also/Too \\\hline
\end{comps}

\begin{remglyph}
\hooglig{Ponds} \remcom{also} contain \remcom{water} to be classical.\\\hline
\end{remglyph}

\begin{prons}
\pronsrtwocone & \textbf{チ (CHI)} & \textbf{いけ (ike)}\\\hline
\pronsrthreecone & --- & --- \\\hline
\end{prons}

\begin{remprons}
\hooglig{Pond} decorated with \comon{cheap} \comkun{ikebana}.\\\hline
\end{remprons}

%{\middel}
% &  \mbox{()}\\\hline
\begin{somme}
池 & \textit{pond} & いけ \mbox{(ike)}\\\hline
電池 & electric + pond = \textit{battery} & デン{\middel}チ \mbox{(DEN{\middel}CHI)}\\\hline
\end{somme}

\begin{decomp}{3}
池&は&干上がった。\\\hline
いけ&は&ひあがった\\\hline
pond&は&to dry up\\\hline
\multicolumn{3}{|r|}{\textit{The pond has dried up.}}\\\hline
\end{decomp}

\end{multicols}

\pagebreak
%%--------------------------------------------------------------------
%21-------------------------------------------------------------------

\begin{multicols}{2}
\begin{glyph}{Regulation (制)}{regulation}\label{glyph:regulation}
Regulation.  Control.
\end{glyph}
\gindex{Regulation}{制}\strindex{8}{制}

\begin{comps}
\comprtwocone & 𠂉 + 市 + 刂\\\hline
\comprthreecone & Man/Person + City + Knife/Sword\\\hline
\end{comps}

\begin{remglyph}
\hooglig{Regulations} in the \remcom{city} about a \remcom{person's} \remcom{sword}.\\\hline
\end{remglyph}

\begin{prons}
\pronsrtwocone & \textbf{セイ (SEI)}& ---\\\hline
\pronsrthreecone & ---& ---\\\hline
\end{prons}

\begin{remprons}
\hooglig{Regulations} concerning \comon{sailing}.\\\hline
\end{remprons}

%{\middel}
% &  \mbox{()}\\\hline
\begin{somme}
自制 & self + regulation = \textit{self-control} & ジ{\middel}セイ \mbox{(JI{\middel}SEI)}\\\hline
制服 & regulation + clothes = \textit{uniform} & セイ{\middel}フク \mbox{(SEI{\middel}FUKU)}\\\hline
制度 & regulation + degree = \textit{system} & セイ{\middel}ド \mbox{(SEI{\middel}DO)}\\\hline
\end{somme}
\end{multicols}

\begin{decomp}{7}
学校&の&制服&は&全く&時代遅れ&だ。\\\hline
ガッコウ&の&セイフク&は&まったく&ジダイおくれ&だ\\\hline
school&の&uniform&は&really&old fashioned&be / is\\\hline
\multicolumn{7}{|r|}{\textit{School uniforms are just out of fashion.}}\\\hline
\end{decomp}


%%--------------------------------------------------------------------
%22-------------------------------------------------------------------
\begin{multicols}{2}
\begin{glyph}{Transport (運)}{transport}\label{glyph:transport}
Transport.\\\hline
A second important meaning (as can be seen in the final compound below) relates to luck.
\end{glyph}
\gindex{Transport}{運}\strindex{12}{運}

\begin{comps}
\comprtwocone & 軍 + 辶 \\\hline
\comprthreecone & Army (\ref{glyph:army}) + Road/Walk \\\hline
\end{comps}

\begin{remglyph}
\hooglig{Transport} on the \remcom{road} to the \remcom{army} base.\\\hline
\end{remglyph}

\begin{prons}
\pronsrtwocone & \textbf{ウン (UN)} & \textbf{はこ (hako)}\\\hline
\pronsrthreecone & --- & --- \\\hline
\end{prons}

\begin{remprons}
\hooglig{Transported} by the \comkun{hakomi} method to \comon{prune} my brain patterns.\\\hline
{\warn}\textit{The Hakomi Method} is a form of mindfulness psychotherapy developed in the 1970s.\\\hline
\end{remprons}

%{\middel}
% &  \mbox{()}\\\hline
\begin{somme}
運ぶ & \textit{transport; carry} & はこ{\middel}ぶ \mbox{(hako{\middel}bu)}\\\hline
運送 & transport + send = \textit{transportation; shipping} & ウン{\middel}ソウ \mbox{(UN{\middel}S\={O})}\\\hline
不運 & not + transport (luck) = \textit{bad luck; misfortune} & フ{\middel}ウン \mbox{(FU{\middel}UN)}\\\hline
\end{somme}

\begin{decomp}{6}
食器&を&台所&に&運んで&ね。\\\hline
ショッキ&を&だいどころ&に&はこんで&ね\\\hline
tableware&を&kitchen&に&to carry&ね\\\hline
\multicolumn{6}{|r|}{\textit{Please carry your dishes to the kitchen.}}\\\hline
\end{decomp}

\end{multicols}

\pagebreak
%%--------------------------------------------------------------------
%23-------------------------------------------------------------------

\begin{multicols}{2}
\begin{glyph}{Mechanism (機)}{mechanism}\label{glyph:mechanism}
The general sense is of mechanisms and things employing a mechanistic action (the final compound below provides an interesting example).\\\hline
This 漢字 also conveys the idea of \textit{opportunity} or \textit{chance} in some of the compounds in which it appears (as in the first compound).\\\hline
This is, incidentally, the only time in the general-use 漢字 that the character for ``how many'' figures as a component in another character.
\end{glyph}
\gindex{Mechanism}{機}\strindex{16}{機}

\begin{comps}
\comprtwocone & 木 + 幾 \\\hline
\comprthreecone & Tree + How Many (\ref{glyph:how_many}) \\\hline
\end{comps}

\begin{remglyph}
\hooglig{A mechanism} for counting \remcom{how many} \remcom{trees} there are.\\\hline
\end{remglyph}

\begin{prons}
\pronsrtwocone & \textbf{キ (KI)} & --- \\\hline
\pronsrthreecone & --- & はた (hata) \\\hline
\end{prons}

\begin{remprons}
\hooglig{Mechanism} for \comon{killing} \ucomkun{hat accessories}.\\\hline
\end{remprons}

%{\middel}
% &  \mbox{()}\\\hline
\begin{somme}
機会 & mechanism + meet = \textit{chance; opportunity} & キ{\middel}カイ \mbox{(KI{\middel}KAI)}\\\hline
電動機 & electric + move + mechanism = \textit{electric motor} & デン{\middel}ドウ{\middel}キ \mbox{(DEN{\middel}D\={O}{\middel}KI)}\\\hline
有機体 & have + mechanism + body = \textit{organism} & ユウ{\middel}キ{\middel}タイ \mbox{(Y\={U}{\middel}KI{\middel}TAI)}\\\hline
\end{somme}

\begin{decomp}{5}
別&の&機会&を&待て。\\\hline
ベツ&の&キカイ&を&まて\\\hline
separate&の&chance&を&to wait\\\hline
\multicolumn{5}{|r|}{\textit{Wait for a second chance.}}\\\hline
\end{decomp}

\end{multicols}

%%--------------------------------------------------------------------
%24-------------------------------------------------------------------
\begin{multicols}{2}
\begin{glyph}{Strong (強)}{strong}\label{glyph:strong}
Strong.
\end{glyph}
\gindex{Strong}{強}\strindex{11}{強}

\begin{comps}
\comprtwocone & 弓 + 厶 + 虫 \\\hline
\comprthreecone & Bow + Self/Private + Insect \\\hline
\end{comps}

\begin{remglyph}
\hooglig{The strong} \remcom{private} crushed the \remcom{insect} with a \remcom{bow}.\\\hline
\end{remglyph}

\begin{prons}
\pronsrtwocone & \textbf{キョウ (KY\={O})} & \textbf{つよ (tsuyo)}\\\hline
\rye{3}{The 訓読み is usually voiced when in the second position.}
\pronsrthreecone & ゴウ (G\={O}) & し (shi) \\\hline
\rye{3}{This 音読み usually appears in words that carry a violent or negative connotation.}
\end{prons}

\begin{remprons}
\hooglig{Strongly} suggest that you \comkun{stew your} \ucomon{gore} during the \comon{Kyoto} \ucomkun{shift}.\\\hline
\end{remprons}

%{\middel}
% &  \mbox{()}\\\hline
\begin{somme}
強い & \textit{strong} & つよ{\middel}い \mbox{(tsuyo{\middel}i)}\\\hline
強化 & strong + change = \textit{strengthening; intensification} & キョウ{\middel}カ \mbox{(KY\={O}{\middel}KA)}\\\hline
強制 & strong + regulation = \textit{compulsion; coercion} & キョウ{\middel}セイ \mbox{(KY\={O}{\middel}SEI)}\\\hline
\end{somme}

\begin{decomp}{5}
風&の&強い&日&です。\\\hline
かぜ&の&つよい&ひ&です\\\hline
wind&の&strong&day&be / is\\\hline
\multicolumn{5}{|r|}{\textit{It's a windy day.}}\\\hline
\end{decomp}

\end{multicols}

\pagebreak
%%--------------------------------------------------------------------
%25-------------------------------------------------------------------

\begin{multicols}{2}
\begin{glyph}{Shine (光)}{shine}\label{glyph:shine}
Shine.  Light.
\end{glyph}
\gindex{Shine}{光}\strindex{6}{光}

\begin{comps}
\comprtwocone & 儿 + 小 + 一 \\\hline
\comprthreecone & Legs/Man/Person + Small + One\\\hline
\end{comps}

\begin{remglyph}
\hooglig{Shine} the light on the \remcom{small}, \remcom{one} \remcom{person}.\\\hline
\end{remglyph}

\begin{prons}
\pronsrtwocone & \textbf{コウ (K\={O})} & \textbf{ひか、ひかり (hika, hikari)}\\\hline
\pronsrthreecone & --- & --- \\\hline
\end{prons}

\begin{remprons}
\hooglig{Shining} example of a \comkun{hick aria} \comon{caused} by a \comkun{hick accent}.\\\hline
{\warn}\textit{A Hick} is an unsophisticated, rural person.\\\hline
\end{remprons}

%{\middel}
% &  \mbox{()}\\\hline
\begin{somme}
光る & \textit{shine} & ひか{\middel}る \mbox{(hika{\middel}ru)}\\\hline
光 & \textit{light; ray} & ひかり \mbox{(hikari)}\\\hline
日光 & sun + shine = \textit{sunshine; sunlight} & ニッ{\middel}コウ \mbox{(NIK{\middel}K\={O})}\\\hline
\end{somme}

\begin{decomp}{5}
月&の&光&が&弱い。\\\hline
つき&の&ひかり&が&よわい\\\hline
moon&の&light&が&weak\\\hline
\multicolumn{5}{|r|}{\textit{The light from the moon is weak.}}\\\hline
\end{decomp}

\end{multicols}

%%--------------------------------------------------------------------
%26-------------------------------------------------------------------
\begin{multicols}{2}
\begin{glyph}{Difficult (難)}{difficult}\label{glyph:difficult}
Difficult.
\end{glyph}
\gindex{Difficult}{難}\strindex{18}{難}

\begin{comps}
\comprtwocone & 艹 + 夫 + 口 + 隹 \\\hline
\comprthreecone & Grass/Herb/Plant + Husband (\ref{glyph:husband}) + Mouth + Old birds \\\hline
\end{comps}

\begin{remglyph}
\hooglig{Difficult} \remcom{husbands} placed \remcom{herbs} in the \remcom{old bird's} \remcom{mouth}.\\\hline
\end{remglyph}

\begin{prons}
\pronsrtwocone & \textbf{ナン (NAN)} & \textbf{むずか、かた (muzuka, kata)}\\\hline
\rye{3}{ナン is by far the most common reading, and the two \mbox{訓読み} form adjectives.\\\hline
When in the second position, かた becomes voiced and makes an adjective of the type shown in the second example below.}
\pronsrthreecone & --- & --- \\\hline
\end{prons}

\begin{remprons}
\hooglig{The difficult} \comon{nun} did not show me the \comkun{moon bazooka} \comkun{catalogue}.\\\hline
\end{remprons}

%{\middel}
% &  \mbox{()}\\\hline
\begin{somme}
難しい & \textit{difficult}  & むずか{\middel}しい \mbox{(muzuka{\middel}shii)}\\\hline
言い難い & say + difficult = \textit{difficult to say}  & い{\middel}い　がた{\middel}い \mbox{(i{\middel}i gata{\middel}i)}\\\hline
難民 & difficult + people = \textit{refugees}  & ナン{\middel}ミン \mbox{(NAN{\middel}MIN)}\\\hline
\end{somme}

\begin{irrr}
有り難う & have + difficult = \textit{Thank you; Thanks} & あ{\middel}り　がと{\middel}う (ari gat\={o})\\\hline
\end{irrr}

\begin{decomp}{4}
勉強&は&難しい&の。\\\hline
ベンキョウ&は&むずかしい&の\\\hline
studies&は&difficult&の\\\hline
\multicolumn{4}{|r|}{\textit{Is the school work hard?}}\\\hline
\end{decomp}

\end{multicols}

\pagebreak
%%--------------------------------------------------------------------
%27-------------------------------------------------------------------

\begin{multicols}{2}
\begin{glyph}{Company (社)}{company}\label{glyph:company}
Company.\\\hline
This 漢字 also refers to \textit{Shinto shrines}; note how the second compound below can encompass both meanings.
\end{glyph}
\gindex{Company}{社}\strindex{7}{社}

\begin{comps}
\comprtwocone & 示礻 + 土 \\\hline
\comprthreecone & Altar/Display + Earth/Soil \\\hline
\end{comps}

\begin{remglyph}
\hooglig{The company's} \remcom{altar} is displayed on holy \remcom{soil}.\\\hline
\end{remglyph}

\begin{prons}
\pronsrtwocone & \textbf{シャ (SHA)} & --- \\\hline
\pronsrthreecone & --- & やしろ (yashiro) \\\hline
\end{prons}

\begin{remprons}
\hooglig{Companies} \comon{shut} their doors because of \comkun{ya sheer obstinacy}.\\\hline
\end{remprons}

%{\middel}
% &  \mbox{()}\\\hline
\begin{somme}
\rye{3}{The final compound below is the only instance in which the 音読み is voiced.}
会社 & meet + company = \textit{company; corporation} & カイ{\middel}シャ \mbox{(KAI{\middel}SHA)}\\\hline
本社 & main + company = \textit{head office; main temple} & ホン{\middel}シャ \mbox{(HON{\middel}SHA)}\\\hline
神社 & god + company = \textit{Shinto Shrine} & ジン{\middel}ジャ \mbox{(JIN{\middel}JA)}\\\hline
\end{somme}

\begin{decomp}{4}
明日&会社&で&ね\\\hline
あした&カイシャ&で&ね\\\hline
tomorrow&company / workplace&で&ね\\\hline
\multicolumn{4}{|r|}{\textit{See you tomorrow in the office.}}\\\hline
\end{decomp}

\end{multicols}

%%--------------------------------------------------------------------
%28-------------------------------------------------------------------
\begin{multicols}{2}
\begin{glyph}{Laugh (笑)}{laugh}\label{glyph:laugh}
Laugh.
\end{glyph}
\gindex{Laugh}{笑}\strindex{10}{笑}

\begin{comps}
\comprtwocone & 竹ケ + 夭 \\\hline
\comprthreecone & Bamboo + Calamity \\\hline
\end{comps}

\begin{remglyph}
\hooglig{Laughing} about the \remcom{bamboo} \remcom{calamity}.\\\hline
\end{remglyph}

\begin{prons}
\pronsrtwocone & \textbf{ショウ (SH\={O})} & \textbf{わら (wara)}\\\hline
\pronsrthreecone & --- & え (e) \\\hline
\end{prons}

\begin{remprons}
\hooglig{Laughter} is \comkun{warranted} at the edge of the \comon{shore}.\\\hline
\end{remprons}

%{\middel}
% &  \mbox{()}\\\hline
\begin{somme}
笑う & \textit{laugh; smile} & さら{\middel}う \mbox{(wara{\middel}u)}\\\hline
高笑い & tall + laugh = \textit{loud laughter} & たか　わら{\middel}い \mbox{(taka wara{\middel}i)}\\\hline
失笑 & lose + laugh = \textit{burst out laughing} & シッ{\middel}ショウ \mbox{(SHIS{\middel}SH\={O})}\\\hline
\end{somme}

\begin{decomp}{3}
皆&が&笑った。\\\hline
みな&が&わらった\\\hline
everyone&が&to laugh\\\hline
\multicolumn{3}{|r|}{\textit{Everybody laughed.}}\\\hline
\end{decomp}

\end{multicols}

\pagebreak
%%--------------------------------------------------------------------
%29-------------------------------------------------------------------

\begin{multicols}{2}
\begin{glyph}{Test (験)}{test}\label{glyph:test}
Test.  Examine.
\end{glyph}
\gindex{Test}{験}\strindex{18}{験}

\begin{comps}
\comprtwocone & 馬 + 央 + 亼\\\hline
\comprthreecone & Horse + Centre + Assembly\\\hline
\end{comps}

\begin{remglyph}
\hooglig{Testing} \remcom{horses} in the \remcom{centre} of the \remcom{assembly}.\\\hline
\end{remglyph}

\begin{prons}
\pronsrtwocone & \textbf{ケン (KEN)} & --- \\\hline
\pronsrthreecone & ゲン (GEN) & --- \\\hline
\end{prons}

\begin{remprons}
\hooglig{Testing} \comon{kennels} made by \ucomon{Genghis} Khan.\\\hline
\end{remprons}

%{\middel}
% &  \mbox{()}\\\hline
\begin{somme}
試験 & try + test = \textit{test; examination} & シ{\middel}ケン \mbox{(SHI{\middel}KEN)}\\\hline
受験 & receive + test = \textit{take an examination} & ジュ{\middel}ケン \mbox{(JU{\middel}KEN)}\\\hline
実験 & actual + test = \textit{experiment} & ジッ{\middel}ケン \mbox{(JIK{\middel}KEN)}\\\hline
\end{somme}

\begin{decomp}{3}
来週&試験&だ。\\\hline
ライシュウ&シケン&だ\\\hline
next week&test / exam&be / is\\\hline
\multicolumn{3}{|r|}{\textit{The exam is coming up next week.}}\\\hline
\end{decomp}

\end{multicols}

%%--------------------------------------------------------------------
%30-------------------------------------------------------------------
\begin{multicols}{2}
\begin{glyph}{History (史)}{history}\label{glyph:history}
History.\\\hline
It is worth taking a moment to see clearly that this 漢字 does \textit{not} function as a component in 更 (``anew'', from Entry \ref{glyph:anew}).
\end{glyph}
\gindex{History}{史}\strindex{5}{史}

\begin{comps}
\comprtwocone & 乂 + 中\\\hline
\comprthreecone & Crossed Swords + Middle\\\hline
\end{comps}

\begin{remglyph}
\hooglig{The history} of the \remcom{Middle East} \remcom{crossing swords}.\\\hline
{\warn}\textit{The Middle East} includes Western Asia, all of Egypt, and Turkey.\\\hline
\end{remglyph}

\begin{prons}
\pronsrtwocone & \textbf{シ (SHI)} & --- \\\hline
\pronsrthreecone & --- & --- \\\hline
\end{prons}

\begin{remprons}
\hooglig{History} is studied to \comon{shield} us from repeating mistakes in the past.\\\hline
\end{remprons}

%{\middel}
% &  \mbox{()}\\\hline
\begin{somme}
\rye{3}{Note that the first two 漢字 in the final compound mean ``literature'' when taken together.}
史学 & history + study = \textit{history; historical studies} & シ{\middel}ガク \mbox{(SHI{\middel}GAKU)}\\\hline
史料 & history + fee (materials) = \textit{historical materials/records} & シ{\middel}リョウ \mbox{(SHI{\middel}RY\={O})}\\\hline
文学史 & literature + history = \textit{literary history} & ブン{\middel}ガク{\middel}シ \mbox{(BUN{\middel}GAKU{\middel}SHI)}\\\hline
\end{somme}
\end{multicols}

\begin{decomp}{8}
学生&の&大&多数&が&史学&を&嫌っている。\\\hline
ガクセイ&の&ダイ&タスウ&が&シガク&を&きらっている\\\hline
students&の&large&majority&が&history&を&to hate / dislike\\\hline
\multicolumn{8}{|r|}{\textit{A majority of students dislike history.}}\\\hline
\end{decomp}
\pagebreak
%%--------------------------------------------------------------------
%2--------------------------------------------------------------------
\vspace*{\fill}
\begin{multicols}{2}
\begin{glyph}{Polish (研)}{polish}\label{glyph:polish}
To polish.  To sharpen.  To whet.
\end{glyph}
\gindex{Polish}{研}\strindex{9}{研}

\begin{comps}
\comprtwocone&石 + 开\\\hline
\comprthreecone&Stone + Start/Initiate\\\hline
\end{comps}

\begin{remglyph}
\hooglig{Polish} the \remcom{stone} from the \remcom{start}.\\\hline
\end{remglyph}

\begin{prons}
\pronsrtwocone&\textbf{ケン (KEN)}&---\\\hline
\pronsrthreecone&---&と (to)\\\hline
\end{prons}

\begin{remprons}
\hooglig{Polish} the \comon{kennel} with a \ucomkun{torrent} of water.\\\hline
\end{remprons}

%{\middel}
% &  \mbox{()}\\\hline
\begin{somme}
{\diff}研ぎ物&polish + thing = \textit{sharpening}&と{\middel}ぎ{\middel}もの \mbox{(to{\middel}gi{\middel}mono)}\\\hline
研究&polish + inquiry = \textit{research}&ケン{\middel}キュウ \mbox{(KEN{\middel}KY\={U})}\\\hline
研究所&research + location = \textit{laboratory; research institute}&ケン{\middel}キュウ{\middel}ジョ \mbox{(KEN{\middel}KY\={U}{\middel}JO)}\\\hline
映研&reflect + polish = \textit{film/cinema studies}&エイ{\middel}ケン \mbox{(EI{\middel}KEN)}\\\hline
研究者&research + individual = \textit{research}&ケン{\middel}キュウ{\middel}シャ \mbox{(KEN{\middel}KY\={U}{\middel}SHA)}\\\hline
研究結果&research + result = \textit{research findings}&ケン{\middel}キュウ{\middel}ケッ{\middel}カ \mbox{(KEN{\middel}KY\={U}{\middel}KEK{\middel}KA)}\\\hline
共同研究&joint + research = \textit{collaborative research}&キョウ{\middel}ドウ{\middel}ケン{\middel}キュウ \mbox{(KY\={O}{\middel}D\={O}{\middel}KEN{\middel}KY\={U})}\\\hline
研学&polish + study = \textit{study}&ケン{\middel}ガク \mbox{(KEN{\middel}GAKU)}\\\hline
通研&passage + polish = \textit{communications laboratory}&ツウ{\middel}ケン \mbox{(TS\={U}{\middel}KEN)}\\\hline
薬研&medicine + polish = \textit{druggist's mortar}&ヤ{\middel}ゲン \mbox{(YA{\middel}GEN)}\\\hline
研究心&research + heart = \textit{(having) an inquiring mind}&ケン{\middel}キュウ{\middel}シン \mbox{(KEN{\middel}KY\={U}{\middel}SHIN)}\\\hline
先行研究&preceding + research = \textit{prior research}&セン{\middel}コウ{\middel}ケン{\middel}キュウ \mbox{(SEN{\middel}K\={O}{\middel}KEN{\middel}KY\={U})}\\\hline
\end{somme}
\end{multicols}

\begin{decomp}{7}
彼&は&日本&文学&の&研究者&だ。\\\hline
かれ&は&ニッポン&ブンガク&の&ケンキュウシャ&だ\\\hline
he&は&Japan&literature&の&researcher&be/is\\\hline
\multicolumn{7}{|r|}{\textit{He's a student of Japanese literature.}}\\\hline
\end{decomp}

\vspace*{\fill}
\pagebreak
%%--------------------------------------------------------------------
